\subsection{Limites et statut épistémologique}

\begin{boitealerte}[Limite~1 --- la couverture n'est pas totale]
La règle se prononce sur $10{,}56\,\%$ des $k$ du corpus. Sur les
$89{,}44\,\%$ restants, le criblage laisse au moins deux survivants :
la fonction renvoie alors « indéterminé ». Le résidu est
\emph{intrinsèque} et non un défaut d'implémentation : sur un
$k$ donné, le nombre de branches survivantes ne dépend que de
$k \bmod M$, et aucune extension finie du crible ne peut faire
disparaître toutes les classes ambguës sans coût prohibitif.
\end{boitealerte}

\begin{boitealerte}[Limite~2 --- la règle n'est pas un test de primalité]
Elle indique \emph{quelle} branche est compatible, non que le candidat
retenu est premier. Pour lever l'ambiguïté résiduelle, un test de
primalité reste nécessaire --- c'est-à-dire un aller-retour vers la
primalité, que la formalisation n'a pas supprimé mais simplement
\emph{repoussé} à un second étage, plus étroit et mieux circonscrit.
\end{boitealerte}

\begin{boitealerte}[Limite~3 --- la partie algébrique conserve un statut
empirique]
La réduction $\SB-\SA = k^{6}B_0$ et la formule fermée de niveau~2 sont
établies et vérifiées. En revanche la \emph{table de correspondance}
entre polarité de la position et signe de la correction, ainsi que
l'ensemble des quatre candidats, sont le fruit d'une exploration
numérique : ils sont \emph{corrects} au sens où ils reproduisent
$676\,191$ ancres sans un seul écart, mais ils ne sont pas
\emph{démontrés} au sens d'un théorème. La section franchit la
frontière entre «~observé~» et «~calculé~» ; elle ne franchit pas celle
entre «~calculé~» et «~prouvé~».
\end{boitealerte}

\begin{boitealerte}[Limite~4 --- la récurrence de la table ne se généralise pas]
Comme établi plus haut, la formule canonique de décomposition du
Digamma est exacte pour $n = 10, 11, 12, 13$ et fausse dès $n = 14$.
Elle est conservée dans le présent document comme \emph{description
de la table d'observation}, à titre d'outil de calcul et de
contrôle, et non comme une loi.
\end{boitealerte}

\begin{boitenote}[Position dans l'architecture générale]
La fonction \textsc{Digamma\_Savard/Cline} intervient exactement là où la
méthode spectrale attendait une information manquante : au niveau~4
(mécanisme d'ancrage), où les quatre essais Digamma étaient jusqu'alors
lancés \emph{simultanément} et sans critère de choix. Elle ne modifie
ni les niveaux~1 à~3 (constantes, formules fermées, construction
terme à terme) ni le niveau~5 (reconstruction pour $n \neq 10$), et ne
modifie pas non plus la formule de reconstruction
\[
  P \;=\; \frac{\SB - \mathrm{Digamma}}{k^{6}},
\]
dont elle ne fournit que le \emph{contenu} de la correction.
\end{boitenote}


\subsection{Implémentation de référence}

\begin{boitenote}[Pseudo-code de la fonction]
\begin{Verbatim}[frame=single,fontsize=\footnotesize]
DIGAMMA_SAVARD(k, n) :
    S_A <- somme_suite_A(k, n)      # arithmétique entière
    S_B <- somme_suite_B(k, n)
    B0  <- k**5 - k**3 + k - 1
    C   <- [B0+k^2, B0-k^2, B0-k, B0+k]  # A(8)-, A(8)+, A(7)+, A(7)-

    survivants <- []
    pour j de 0 a 3 :
        si tous(C[j] mod p != 0 pour p dans [3,5,7,11,13,17,19,23]) :
            ajouter j a survivants

    si longueur(survivants) == 1 :
        j <- survivants[0]
        signe   <- -1 si j dans {0,3} sinon +1
        exposant <- 8 si j dans {0,1} sinon 7
        Digamma <- S_A + signe * k**exposant
        retourner (Digamma, (S_B - Digamma) / k**6)
    sinon :
        retourner (INDETERMINE, branches survivantes)
\end{Verbatim}
\end{boitenote}

\begin{boitenote}[Précautions d'implémentation]
\begin{itemize}[nosep]
  \item \textbf{Arithmétique entière obligatoire.} Pour $k \ge 50$, les
    valeurs de $\SA$ et $\SB$ dépassent $10^{15}$~: l'arithmétique
    flottante introduit des erreurs d'arrondi qui se propagent au
    test de divisibilité et faussent le verdict de criblage.
  \item \textbf{Le test de divisibilité précède toute hypothèse de
    primalité.} La fonction ne doit jamais court-circuiter sur un
    «~candidat préféré~»~: c'est précisément ce court-circuit qui
    réintroduirait la circularité.
  \item \textbf{Le cas indéterminé doit être rendu visible.} En cas
    d'ambiguïté, l'agent doit signaler \textsc{C\_NON\_CERTIFIÉ} et
    attendre une validation externe, conformément au contrat HOL,
    plutôt que de trancher silencieusement.
\end{itemize}
\end{boitenote}


\subsection{Conclusion}

La correction Digamma, longtemps traitée comme une valeur observée,
devient ici le résultat d'une \emph{fonction calculable} : la
\emph{fonction \textsc{Digamma\_Savard/Cline}}. Sa structure est la
suivante.

\begin{enumerate}[nosep,leftmargin=*]
  \item \textbf{Une réduction algébrique.} Les quatre candidats
    possibles se ramènent à quatre polynômes
    $B_0 \pm k^{2}$ et $B_0 \pm k$, où
    $B_0(k) = k^{5}-k^{3}+k-1$.
  \item \textbf{Une règle de sélection par criblage modulaire.} La
    polarité de la position dans la suite A commande le signe de la
    correction ; ce signe est \emph{déduit} par élimination
    modulaire, jamais \emph{choisi}.
  \item \textbf{Une validation rigoureuse.} Zéro violation
    d'intégrité sur $676\,191$ ancres~; exactitude de $100\,\%$ sur
    les $71\,401$ ancres à survivant unique~; effondrement de la
    pureté à $0$--$8{,}6\,\%$ sous permutation des signes~; stabilité
    de l'exactitude à $100\,\%$ lorsque la couverture passe de
    $10{,}56\,\%$ à $49{,}65\,\%$.
  \item \textbf{Un statut honnête.} La règle est exacte là où elle se
    prononce et le dit quand elle ne peut pas se prononcer~; elle ne
    prétend ni tester la primalité, ni démontrer la table de
    correspondance.
\end{enumerate}

\begin{boiteformule}[Énoncé final]
Pour tout $k \ge 2$, la fonction \textsc{Digamma\_Savard/Cline} retourne
une correction calculée, non observée, dont le signe est déterminé par
la seule arithmétique modulaire~; et elle retourne «~indéterminé~»
plutôt que de trancher au hasard. \emph{La fonction ne prétend pas
résoudre la primalité~: elle déplace le problème du \emph{choix} vers
la \emph{contamination}.}
\end{boiteformule}
